\documentclass[]{article}

\usepackage{amsmath}
\usepackage{multirow}
\usepackage{natbib}
\usepackage{graphicx} 


\begin{document}

The results of our investigation are presented in two parts. First, we quantify the computational performance of our GPU-accelerated simulation. Second, we present a detailed physical validation of the simulation's accuracy by analyzing the ensemble-averaged matter power spectrum.

\subsection{Computational performance}
A primary objective of this work was to leverage GPU acceleration to drastically reduce the computational time required for cosmological N-body simulations. Our implementation using the NVIDIA Warp framework demonstrates a significant performance improvement over traditional CPU-based methods.

Each simulation in our 10-realization ensemble, evolving $512^3$ particles from $z=127$ to $z=0$ over 200 time steps, completed in approximately 20 seconds on a single GPU. This represents a substantial reduction from the estimated wall-clock time of approximately 5 hours for a comparable PM simulation running on a multi-core CPU. The total computational time for the entire ensemble was approximately 200 seconds. This dramatic increase in speed enables the rapid generation of large simulation suites, which is critical for modern cosmological analyses that require robust statistical sampling and covariance estimation.

\subsection{Physical validation}
To assess the physical accuracy of our simulation, we computed the matter power spectrum, $P(k)$, for each of the ten independent realizations at $z=0$. These simulations were initialized using second-order Lagrangian Perturbation Theory (2LPT) to ensure accurate initial particle displacements and velocities. The ensemble-averaged power spectrum, $\langle P_{\text{sim}}(k) \rangle$, provides a statistically robust measurement by mitigating the effects of cosmic variance.

Figure \ref{fig:pk_ensemble_mean} compares the ensemble-averaged power spectrum from our simulations to the non-linear theoretical prediction from CAMB HaloFit. The ratio plot reveals three distinct regimes of agreement. On large cosmological scales ($k < 0.03 \, \text{h/Mpc}$), our simulation accurately reproduces the theoretical power spectrum, with the ratio being approximately 0.96, achieving our target of 5\% agreement. At intermediate scales ($0.03 \le k \le 0.1 \, \text{h/Mpc}$), we observe a systematic underprediction of power on the order of 10-15\%. At smaller scales ($k > 0.1 \, \text{h/Mpc}$), the deficit in power becomes more pronounced. This behavior is an expected and well-understood limitation of the Particle-Mesh (PM) method, where the gravitational force resolution is determined by the grid spacing. The finite grid acts as a low-pass filter, smoothing out density fluctuations and suppressing power near the Nyquist frequency of the grid.

\begin{figure}[h!]
    \centering
    \includegraphics[width=0.5\textwidth]{../input_files/plots/step_6_plot_0_1778590772.png}
    \caption{Ensemble-averaged matter power spectrum from 10 simulations using 2LPT initial conditions. Top: The mean simulated power spectrum, $\langle P_{\text{sim}}(k) \rangle$, is compared to the non-linear prediction from CAMB HaloFit, $P_{\text{CAMB, NL}}(k)$. Bottom: The ratio of the simulated power to the theoretical prediction. The simulation agrees with the theoretical model to within 5\% at large scales ($k < 0.03 \, \text{h/Mpc}$), while systematically underpredicting power at smaller scales ($k > 0.1 \, \text{h/Mpc}$).}
    \label{fig:pk_ensemble_mean}
\end{figure}

The use of an ensemble is critical for obtaining a robust measurement of the mean power, particularly at large scales where there are few independent modes in the simulation volume. Figure \ref{fig:pk_ensemble_full} illustrates this by showing the same mean power spectrum from Figure \ref{fig:pk_ensemble_mean} along with the $1\sigma$ standard deviation across the 10 realizations, which quantifies the cosmic variance. As shown, the cosmic variance is substantial at the largest scales ($k \le 0.05 \, \text{h/Mpc}$), where the standard deviation can exceed 20\% of the mean power. At smaller scales ($k \ge 0.1 \, \text{h/Mpc}$), the cosmic variance becomes negligible ($<2\%$), and the ensemble average converges to a stable measurement of the systematic deviation from the theoretical model.

\begin{figure}[h!]
    \centering
    \includegraphics[width=0.5\textwidth]{../input_files/plots/pk_ensemble.png}
    \caption{The ensemble-averaged matter power spectrum from Figure \ref{fig:pk_ensemble_mean} shown with the $1\sigma$ cosmic variance across the 10 realizations (shaded region). The variance is significant at large scales ($k < 0.05 \, \text{h/Mpc}$), highlighting the need for an ensemble to obtain a robust measurement of the mean. At smaller scales, the variance becomes negligible compared to the systematic deviation from the theoretical model.}
    \label{fig:pk_ensemble_full}
\end{figure}

These results confirm that our GPU-accelerated simulation provides reliable predictions on the large scales ($k < 0.1 \, \text{h/Mpc}$) most relevant for many cosmological analyses, such as studies of Baryon Acoustic Oscillations, while accurately capturing the expected limitations of the PM method at smaller scales.

\end{document}
                